\documentclass[11pt]{article}
\usepackage[a4paper,margin=23mm]{geometry}
\usepackage{amsmath,amssymb,amsthm,mathtools,microtype}
\usepackage[hidelinks]{hyperref}
\newtheorem{theorem}{Theorem}
\newtheorem{lemma}{Lemma}
\newcommand{\Q}{\mathcal Q}
\newcommand{\W}{\mathcal W}
\newcommand{\tr}{\operatorname{tr}}
\newcommand{\sxdeg}{\operatorname{sxdeg}}
\setlength{\parskip}{2pt}
\title{An exact planar quantum region beyond finite NPA depth}
\author{}
\date{27 October 2026}
\begin{document}
\maketitle
\vspace{-2em}

In the simplest Bell experiment, retain two observed quantities:
\[
 M=\langle A_0\rangle+\langle B_0\rangle,\qquad
 S=\langle A_0B_0+A_0B_1+A_1B_0-A_1B_1\rangle .
\]
Their quantum region is
\(\W=\{(M,S):p\in\Q_{222}\}\).
The supplied symmetric-tilt construction describes upper bounds on \(S+aM\).
We extend it to every linear functional \(hS+aM\), including negative \(h\),
and obtain an explicit membership test for the whole region \(\W\).

\begin{theorem}
The region \(\W\) is the projection of one fixed affine semidefinite set with
PSD blocks of orders \(5\) and \(4\), given explicitly in
Eqs.~\eqref{eq:dualmat}--\eqref{eq:coords}.
Its semidefinite extension degree is exactly two.
For the same region, every fixed standard NPA level has divergent relative
support error along the symmetric doubly tilted family as \(a\uparrow1\),
by the prior nontermination theorem~\cite{prior}.
\end{theorem}

The explicit blocks and the extension-degree assertion serve different purposes.
The former give a reproducible lift; the latter establish the smallest possible
largest block order among all finite lifts, without specifying the number of blocks.
The proof reduces all measurement pairs to equal overlaps, converts the remaining
matrix polynomial into a sum of squares, and dualizes its affine coefficient equations.

\section{One polynomial for every symmetric support direction}

Let \(H=h\,\mathrm{CHSH}+a(A_0+B_0)\), with \(h,a\in\mathbb R\).
Write
\[
 {\cal C}=\{(\lambda,a,h):\lambda I-H\succeq0
 \text{ for every quantum realization}\}.
\]
Jordan's lemma, after projective dilation, reduces this condition to real
two-qubit blocks. Put
\[
 A_0=B_0=Z,\quad A_1=cZ+rX,\quad B_1=dZ+sX,\qquad
 r^2=1-c^2,\ s^2=1-d^2 .
\]
We may take \(r,s\ge0\) by local \(Z\) conjugations. Endpoints include the
one-dimensional blocks. Complex state vectors cause no further condition,
because the Bell matrix is real symmetric.

Set \(H_*=|h|\), \(b=|a|\) and \(L=2H_*+2b\).
Deterministic assignments attain \(L\): choose \(A_0=B_0=\operatorname{sign}(a)\),
and take \(A_1=B_1\) equal to that sign for \(h\ge0\), or its negative for \(h<0\).
At \(a=0\), either sign works. Therefore every admissible \(\lambda\) satisfies
\(\lambda\ge L\).

\paragraph{The determinant controls positivity.}
With \(\Omega=1+c+d-cd\in[-2,2]\), the coefficients in
\[
 \det(uI+\lambda I-H)=u^4+e_1u^3+e_2u^2+e_3u+e_4
\]
are
\begin{align*}
 e_1&=4\lambda,&
 e_2&=6\lambda^2-4a^2-8h^2,\\
 e_3&=4\lambda^3-8\lambda a^2-16\lambda h^2-8ha^2\Omega,
 &e_4&=\det(\lambda I-H).
\end{align*}
The first three are nonnegative for \(\lambda\ge L\).
For \(e_3\), bound its last term below by \(-16H_*b^2\).
The resulting cubic is increasing for \(\lambda\ge L\), and at \(L\) equals
\(16b(b+2H_*)^2\).
Consequently
\[
 \lambda I-H\succeq0\quad\Longleftrightarrow\quad e_4\ge0 .
\]
Indeed, a negative eigenvalue would produce a positive root of a polynomial
whose coefficients are all nonnegative. This argument includes zero determinants
and multiple zero eigenvalues.

\paragraph{Equal overlaps give the global minimum.}
Put \(x=(1-c)/2\), \(y=(1-d)/2\), and denote the determinant by \(D(x,y)\).
Direct expansion gives
\begin{align*}
D(x,y)={}&D_0+64h^2(4h^2-a^2)xy(x+y)\\
 &+(32\lambda ha^2+128h^2a^2-256h^4)xy-256h^4x^2y^2,\\
D_0={}&[(\lambda-2h)^2-4a^2](\lambda+2h)^2.
\end{align*}
Hence
\begin{equation}
 D(x,y)-D(\sqrt{xy},\sqrt{xy})
 =64h^2(4h^2-a^2)xy(\sqrt x-\sqrt y)^2 .
 \label{eq:diagonal}
\end{equation}
For \(|a|\le |h|\), the right side is nonnegative. Testing all equal-overlap
blocks therefore suffices for every pair \(c,d\). No division by an angle,
a determinant, or \(h\) occurs.

\paragraph{The remaining directions are classical.}
The range \(b\ge H_*\) has exact support value \(L\).
To prove this, set \(C_B=(B_0+B_1)/2\), \(D_B=(B_0-B_1)/2\). Then
\begin{equation}
 4I-(\mathrm{CHSH}+2A_0)
 =(I-C_B-A_1D_B)^2+
 \frac14\sum_{y=0}^1[(I-A_0)(I+B_y)]^2\succeq0 .
 \label{eq:critical}
\end{equation}
Here \(C_B^2+D_B^2=I\) and \(C_BD_B+D_BC_B=0\).
Party exchange and output flips give
\(\sigma\mathrm{CHSH}+2\tau A_0\le4I\) and
\(\sigma\mathrm{CHSH}+2\tau B_0\le4I\), for either sign \(\sigma,\tau\).
Average these inequalities, multiply by \(H_*\), and add
\((b-H_*)[2I-\tau(A_0+B_0)]\succeq0\).
The resulting bound is \(H\le LI\), attained by the deterministic strategies above.
This also covers \(h=0\), the critical line, and both coefficient signs.

\paragraph{The quartic matrix.}
On \(c=d\), the Bell basis
\[
 \phi_\pm=(|00\rangle\pm|11\rangle)/\sqrt2,\qquad
 \psi_\pm=(|01\rangle\pm|10\rangle)/\sqrt2
\]
splits off \(\psi_-\), with eigenvalue \(-2hc\).
In the basis \((\phi_+,\phi_-,\psi_+)\), the remaining block is
\[
 H_3(c)=
 \begin{pmatrix}
 2hc&2a&0\\
 2a&2h(1+c-c^2)&2h(1-c)\sqrt{1-c^2}\\
 0&2h(1-c)\sqrt{1-c^2}&-2h(1+c-c^2)
 \end{pmatrix}.
\]
Use \(c=(1-t^2)/(1+t^2)\) and multiply its slack by \((1+t^2)^2\).
Writing
\[
 p_0=(1+t^2)^2,\qquad p_1=1-t^4,\qquad p_2=1+4t^2-t^4,
\]
the resulting quartic is
\begin{equation}
 K_{h,a,\lambda}(t)=
 \begin{pmatrix}
 \lambda p_0-2hp_1&-2ap_0&0\\
 -2ap_0&\lambda p_0-2hp_2&-8ht^3\\
 0&-8ht^3&\lambda p_0+2hp_2
 \end{pmatrix}.
 \label{eq:K}
\end{equation}
Finite \(t\ge0\) covers \(-1<c\le1\); \(c=-1\) follows by continuity.
Negative \(t\) adds no condition because
\(K(-t)=\operatorname{diag}(1,1,-1)K(t)\operatorname{diag}(1,1,-1)\).

Conversely, \(K(t)\succeq0\) for all real \(t\) forces the local bound.
Its upper-left block at \(t=0\) gives \(\lambda\ge2h+2|a|\);
the leading coefficient at \(t=\infty\) gives \(\lambda\ge-2h+2|a|\).
Thus \(\lambda\ge L\), and the omitted scalar sector is controlled.
Equation~\eqref{eq:diagonal} handles the low-tilt range; Eq.~\eqref{eq:critical}
handles the rest. We have proved, for all real coefficients,
\begin{equation}
 (\lambda,a,h)\in{\cal C}
 \quad\Longleftrightarrow\quad
 K_{h,a,\lambda}(t)\succeq0\quad\text{for every }t\in\mathbb R .
 \label{eq:cone}
\end{equation}

\paragraph{A fixed affine Gram lift.}
Every PSD real univariate matrix polynomial is a polynomial matrix sum of squares,
including identically singular ones~\cite{hs}.
Since the diagonal entries of \(K\) have degree at most four, every square factor
has degree at most two: highest-degree real squares cannot cancel.
Average a Gram representation under the involution above. Its two parity blocks
use precisely
\[
 Z_+=(\xi_1,\xi_2,t\xi_3,t^2\xi_1,t^2\xi_2)^T,\qquad
 Z_-=(\xi_3,t\xi_1,t\xi_2,t^2\xi_3)^T.
\]
Thus \(\xi^TK(t)\xi=Z_+^TGZ_++Z_-^TJZ_-\), with
\(G\succeq0\) of order five and \(J\succeq0\) of order four.
The sixteen coefficient equations are
\begin{equation}
\begin{aligned}
 g_{11}&=\lambda-2h,&g_{12}&=-2a,\\
 g_{22}&=\lambda-2h,&j_{11}&=\lambda+2h,\\
 g_{13}+j_{12}&=0,&g_{23}+j_{13}&=0,\\
 2g_{14}+j_{22}&=2\lambda,&g_{15}+g_{24}+j_{23}&=-4a,\\
 2g_{25}+j_{33}&=2\lambda-8h,&g_{33}+2j_{14}&=2\lambda+8h,\\
 g_{34}+j_{24}&=0,&g_{35}+j_{34}&=-8h,\\
 g_{44}&=\lambda+2h,&g_{45}&=-2a,\\
 g_{55}&=\lambda+2h,&j_{44}&=\lambda-2h .
\end{aligned}
\label{eq:rows}
\end{equation}
These equations include zero coefficients, lower polynomial degrees and every
rank-deficient case. Setting \(h=1\) gives the common-tilt epigraph for all real
\(a\); its restriction to \(0\le a\le1\) additionally requires those two scalar
inequalities.

\section{Explicit membership matrices for the quantum region}

Introduce sixteen real variables \(z_1,\ldots,z_{16}\) and define
\begin{equation}
\begin{split}
 U(z)&=
 \begin{pmatrix}
 z_1&z_2/2&z_5/2&z_7&z_8/2\\
 z_2/2&z_3&z_6/2&z_8/2&z_9\\
 z_5/2&z_6/2&z_{10}&z_{11}/2&z_{12}/2\\
 z_7&z_8/2&z_{11}/2&z_{13}&z_{14}/2\\
 z_8/2&z_9&z_{12}/2&z_{14}/2&z_{15}
 \end{pmatrix},\\[2pt]
 V(z)&=
 \begin{pmatrix}
 z_4&z_5/2&z_6/2&z_{10}\\
 z_5/2&z_7&z_8/2&z_{11}/2\\
 z_6/2&z_8/2&z_9&z_{12}/2\\
 z_{10}&z_{11}/2&z_{12}/2&z_{16}
 \end{pmatrix}.
\end{split}\label{eq:dualmat}
\end{equation}
The exact membership test is
\[
 (M,S)\in\W\quad\Longleftrightarrow\quad
 \text{there exists }z\in\mathbb R^{16}
 \text{ with }U(z)\succeq0,\quad V(z)\succeq0,
\]
and
\begin{align}
 1={}&z_1+z_3+z_4+2z_7+2z_9+2z_{10}+z_{13}+z_{15}+z_{16},\nonumber\\
 M={}&2z_2+4z_8+2z_{14},\nonumber\\
 S={}&2z_1+2z_3-2z_4+8z_9-8z_{10}+8z_{12}
       -2z_{13}-2z_{15}+2z_{16}.
 \label{eq:coords}
\end{align}

Here is why dualization gives both inclusions and attains the auxiliary variables.
Write Eq.~\eqref{eq:rows} as
\(\mathcal A(G,J)=B(\lambda,a,h)\), in the displayed row order.
The matrices in Eq.~\eqref{eq:dualmat} are exactly
\(\mathcal A^*z\), and
\(B^Tz=(\ell(z),-M(z),-S(z))\), where \(\ell\) is the normalization in
Eq.~\eqref{eq:coords}. Therefore a feasible \(z\) pairs nonnegatively with every
valid support inequality.

For the reverse implication, the dual image
\[
 {\cal D}=\{B^Tz:U(z)\succeq0,\ V(z)\succeq0\}
\]
is closed. Indeed, the positive definite matrices
\[
 G_0=\operatorname{diag}(1,1,2,1,1),\qquad
 J_0=\operatorname{diag}(1,2,2,1)
\]
satisfy \(\mathcal A(G_0,J_0)=B(1,0,0)\), and hence
\(\ell(z)=\tr[U(z)G_0]+\tr[V(z)J_0]\).
If \(B^Tz\) converges, this weighted trace bounds both PSD matrices.
Every \(z_i\) occurs in their entries, so \(z\) is bounded and has a convergent
subsequence. The limit is feasible. The strictly feasible pair \((G_0,J_0)\)
licenses the conic intersection-dual rule, giving
\({\cal C}^*=\overline{\cal D}={\cal D}\).
Finally, the compact convex set \(\W\) is the intersection of its support
half-spaces, so \((M,S)\in\W\) exactly when \((1,-M,-S)\in{\cal C}^*\).
This proves the membership test. The normalized feasible \(z\)-set is itself
compact by the same trace argument.

\section{Minimal block order and the NPA comparison}

The set \(\W\) is compact and convex. The explicit lift makes it semialgebraic.
Scheiderer's planar theorem therefore gives \(\sxdeg(\W)\le2\)~\cite{soc}.
A short argument shows that \(\W\) is not a polygon.

Let \(f(a)=\max_{\W}(S+aM)\), and \(w=2\sqrt2=f(0)\).
The identity
\[
 \mathrm{CHSH}(c,d)^2=4I+4rs\,Y\otimes Y
\]
shows that attaining \(w\) requires \(r=s=1\).
At \(c=d=0\), the top state is
\(v=\cos(\pi/8)\phi_-+\sin(\pi/8)\psi_+\), with both reduced states \(I/2\).
Thus every CHSH-maximizing behavior has \(M=0\).
For maximizers \(p_a\) of \(f(a)\), compactness then gives \(M(p_a)\to0\) as
\(a\downarrow0\), and
\[
 0\le\frac{f(a)-w}{a}\le M(p_a)\longrightarrow0.
\]
Yet \(f(a)>w\) for every \(a>0\): keeping \(c=d=0\), the normalized state
\((v+\varepsilon\phi_+)/\sqrt{1+\varepsilon^2}\) gives
\[
 \frac{w+4a\cos(\pi/8)\varepsilon}{1+\varepsilon^2}>w
 \quad\text{for}\quad \varepsilon=\frac{2a\cos(\pi/8)}w .
\]
A polygon would make \(f\) piecewise affine with finitely many pieces.
Its zero right derivative at zero would force an initial constant piece,
contradicting this strict increase. Therefore \(\W\) is nonpolyhedral, and
\[
 \boxed{\sxdeg(\W)=2.}
\]
The same proof gives extension degree two for
\(\{(a,\lambda):0\le a\le1,\ \lambda\ge f(a)\}\).

The prior NPA result applies to these very support directions \(h=1\)~\cite{prior}:
\[
 f(1-T)=4-2T+\frac43T^3+O(T^4),\qquad
 \frac{\omega_{\mathrm{NPA}}^{(L)}(1-T)-f(1-T)}{T^3}\longrightarrow+\infty
\]
for every fixed finite level \(L\).
Consequently no fixed NPA projection onto \((M,S)\) equals \(\W\), while
Eqs.~\eqref{eq:dualmat}--\eqref{eq:coords} describe that same region exactly.
The nontermination theorem is an imported established result; the finite lift
and nonpolyhedrality arguments above do not depend on its proof.

\section{All signed asymmetric tilt branches}

The weighted-ray argument also extends to arbitrary real tilts in
\[
 H_{\alpha,\beta}=\mathrm{CHSH}+\alpha A_0+\beta B_0.
\]
Put \(q=|\alpha|+|\beta|\). Its local bound is \(2+q\), attained for any tilt
signs. If \(q\ge2\), signed one-sided versions of Eq.~\eqref{eq:critical},
convex combination on \(q=2\), and sign-adjusted marginal slacks above it prove
that this is also the quantum bound.

Suppose \(q<2\) and \(\lambda\ge2+q\). For the full Jordan slack,
\[
 e_2=6\lambda^2-2\alpha^2-2\beta^2-8,\qquad
 e_3=4\lambda^3-4\lambda(\alpha^2+\beta^2+4)-8\alpha\beta\Omega .
\]
Both are nonnegative. For \(e_3\), a lower bound is
\(g(\lambda)=4\lambda^3-4\lambda(q^2+4)-4q^2\).
It increases on this range, and \(g(2+q)=32q+12q^2\ge0\).
Thus the determinant again decides PSD.

With \(x=(1-c)/2\), \(y=(1-d)/2\), direct expansion gives
\begin{align*}
 D(x,y)&=D_0+xy[C+64Ax+64By-256xy],\\
 A&=4-\alpha^2,\qquad B=4-\beta^2,\\
 C&=32\lambda\alpha\beta+64\alpha^2+64\beta^2-256,\\
 D_0&=[(\lambda-2)^2-(\alpha+\beta)^2]
       [(\lambda+2)^2-(\alpha-\beta)^2].
\end{align*}
Both \(A,B\) are positive. Exchange parties if necessary so that
\(|\alpha|\ge|\beta|\), and set \(\rho=A/B\in(0,1]\).
For a fixed product \(p=xy\), minimize \(Ax+Bp/x\) on \(p\le x\le1\).
For \(p\le\rho\), its minimizer lies on \(x=z,y=\rho z\).
For \(p>\rho\), it lies on \(x=1\). On that edge,
\[
 D(1,y)=-F_-(y)F_+(y),
\]
where
\[
 F_-=-(\lambda+\alpha)^2+\beta^2+4+4\beta(2y-1),\qquad
 F_+=(\lambda-\alpha)^2-\beta^2-4+4\beta(2y-1).
\]
Since \(\lambda\pm\alpha\ge2+|\beta|\), we have \(F_-\le0\le F_+\).
Hence the complete measurement square reduces to the weighted ray, for every
choice of tilt signs.

Along the ray, put
\(R=4z(1-z)\), \(T=4\rho z(1-\rho z)\), and \(\Omega=2-4\rho z^2\).
Congruence by \(\operatorname{diag}(\sqrt{RT},\sqrt R,\sqrt T,1)\) gives
\begin{equation}
 P(z)=
 \begin{pmatrix}
 RT(\lambda-\Omega-\alpha-\beta)&-2zRT&-2\rho zRT&RT\\
 -2zRT&R(\lambda+\Omega-\alpha+\beta)&RT&2\rho zR\\
 -2\rho zRT&RT&T(\lambda+\Omega+\alpha-\beta)&2zT\\
 RT&2\rho zR&2zT&\lambda-\Omega+\alpha+\beta
 \end{pmatrix}.
 \label{eq:rayP}
\end{equation}
This matrix has degree at most six. The congruence is invertible on \(0<z<1\);
the original slack at both endpoints follows by continuity from that interval.
The matrix interval sum-of-squares theorem gives
\begin{equation}
 P(z)=Y_3(z)^TUY_3(z)+z(1-z)Y_2(z)^TVY_2(z),\qquad
 U\succeq0,\quad V\succeq0,
 \label{eq:interval}
\end{equation}
with \(Y_j=(I_4,zI_4,\ldots,z^jI_4)^T\), and block orders \(16,12\).
One can obtain this interval theorem directly from the real-line matrix theorem:
apply it to \((1+t^2)^6P(t^2/(1+t^2))\), split its factors into even and odd parts,
and divide back. Even factors become polynomials of degree at most three in \(z\);
odd factors give \(\sqrt{z(1-z)}\) times polynomials of degree at most two.
This argument also includes singular matrices and endpoint zeros.

For fixed \(\alpha,\beta\), coefficient matching in Eq.~\eqref{eq:interval}
is affine in \(\lambda,U,V\). Include \(\lambda\ge2+q\) explicitly.
This yields an exact fixed-size SDP throughout the signed subcritical diamond,
including either zero-tilt axis and CHSH. The critical and exterior branches
are the elementary local-bound certificates already proved.
The ratio \(\rho\) still depends nonlinearly on the two independent tilts:
this construction computes their support values, while the common-tilt
construction above gives an affine membership lift.

\section*{How the results fit together}

Three general steps organize the proofs. First, a local deterministic bound makes
the first three characteristic coefficients positive, leaving one determinant.
Second, minimizing that determinant at fixed angle product removes one testing
variable. Third, univariate matrix positivity gives bounded-degree square factors;
when its coefficients remain affine in the visible Bell parameters, duality gives
a single membership lift.

For equal tilts these steps close the entire planar region, including the NPA-hard
directions. For independent tilts they close every fixed-input support calculation.
The existing five-parameter C1 neighborhood uses a complementary mechanism:
positive Grams on contact complements persist under compatible changes of the
source coefficients. Its eighteen-function space remains fixed while the contact
moves. It supplies local source coverage beyond the exact C1 slice.

The remaining global construction has two substantive admission targets:
every strict C1 source and every strict C2 source must enter one fixed sound
finite library, or an adequate replacement. The existing reflection, singular
coefficient density, event and closed-face cones then supply the other branches.
The planar theorem above adds an exact region and a completed NPA comparison;
the broader branch map records where the universal source obligations remain.

\begin{thebibliography}{9}
\bibitem{hs} C. Hanselka and R. Sinn, Positive semidefinite univariate matrix
polynomials, \emph{Math. Z.} \textbf{292}, 83--101 (2019);
\href{https://arxiv.org/pdf/1707.08261}{Corollary 2.8}.
\bibitem{soc} C. Scheiderer, Second-Order Cone Representation for Convex Sets in
the Plane, \emph{SIAM J. Appl. Algebra Geom.} \textbf{5}, 114--139 (2021),
\href{https://doi.org/10.1137/20M133717X}{doi:10.1137/20M133717X};
\href{https://arxiv.org/html/2004.04196v2}{Theorem 0.1}.
\bibitem{prior} A. Chaturvedi, \emph{No Finite NPA Level Characterizes the Complete
Quantum Set in the Simplest Bell Scenario}, prior Letter and Supplemental
Material supplied in the project.
\end{thebibliography}
\end{document}
